\begin{testscenario}{27}{3}
  \tcmeta{TC27}{CS 2568/01}{ระบบเว็บแอปพลิเคชัน ``NUCARE''}{นายสมชาย ใจดี}{วันที่ 5 พฤศจิกายน พ.ศ.\ 2568}{1.0}
  \tcname{การแบ่งหน้ารายการยาวแบบเดียวกับหน้า admin}
  \tcpre{มีชุด helper \tccode{list-pagination} และหน้าที่มีรายการจำนวนมาก เช่น ประวัติคำขอปรึกษาของนิสิต / แท็บปิดแล้วของที่ปรึกษา}
  \tcdesc{ตรวจสอบว่าการตัดหน้าและคำนวณจำนวนแถวต่อหน้าทำงานถูกต้อง ไม่เทรายการทั้งหมดลงหน้าจอ}

  \tccase{1}{ตัดหน้ารายการและแสดงช่วง X--Y จาก Z}{%
    \begin{tcsteps}
      \item เรียก \tccode{slicePageItems} กับรายการ 25 ชิ้น ขนาดหน้า 5
      \item ตรวจหน้า 1 และหน้าสุดท้าย
      \item ตรวจกรณีหน้าเกินจำนวนหน้าและรายการว่าง
    \end{tcsteps}}{คืนช่วงรายการและ \tccode{totalPages} ถูกต้อง พร้อม clamp หน้า}{รัน \tccode{npx tsx --test src/lib/list-pagination.test.ts} ผ่านทั้งชุด \tccode{slicePageItems}}{Pass}

  \tccase{2}{คำนวณแถวต่อหน้าตามขนาดจอ}{%
    \begin{tcsteps}
      \item เรียก \tccode{computeAdaptiveRowsPerPage} สำหรับเดสก์ท็อปและมือถือ
      \item ตรวจค่าอยู่ในช่วง min--max และไม่เกิน max เมื่อจอสูงมาก
    \end{tcsteps}}{จำนวนแถวอยู่ในช่วงที่กำหนดตาม viewport}{เคสในชุด \tccode{computeAdaptiveRowsPerPage} ผ่านครบ}{Pass}

  \tccase{3}{UI รายการยาวใช้ footer ก่อนหน้า/ถัดไป}{%
    \begin{tcsteps}
      \item เปิด \tccode{/student/advisor} ด้วยบัญชีที่มีประวัติจำนวนมาก
      \item ตรวจก้อนประวัติคำขอปรึกษาว่ามีข้อความแสดงช่วงรายการและปุ่มแบ่งหน้า
      \item เปิดคิวคำขอปรึกษาฝั่งที่ปรึกษาแท็บปิดแล้วแล้วสลับหน้า
    \end{tcsteps}}{ไม่เทรายการทั้งหมดในครั้งเดียว และเปลี่ยนหน้าได้}{หน้านิสิตและที่ปรึกษาใช้ \tccode{ListPaginationFooter} เมื่อจำนวนเกินขนาดหน้า}{Pass}
\end{testscenario}
